\documentclass[11pt,a4paper]{article}
\usepackage[margin=2.4cm]{geometry}
\usepackage{amsmath,amssymb}
\usepackage{booktabs,longtable,array,tabularx}
\usepackage[dvipsnames]{xcolor}
\usepackage[most]{tcolorbox}
\usepackage{listings}
\usepackage{graphicx}
\usepackage{float}
\usepackage{caption}
\usepackage{enumitem}
\usepackage[expansion=false]{microtype}
\usepackage{fancyhdr}
\usepackage{xurl}
\usepackage[colorlinks=true,linkcolor=NavyBlue,urlcolor=NavyBlue,citecolor=NavyBlue]{hyperref}

\definecolor{navy}{HTML}{14213D}
\definecolor{crim}{HTML}{C0392B}
\definecolor{teal}{HTML}{1F8A7A}
\definecolor{paper}{HTML}{EEF2F8}
\definecolor{gold}{HTML}{FFF4DC}

% ---- boxes -------------------------------------------------------------
\newtcolorbox{defbox}[1]{breakable,colback=paper,colframe=navy,fonttitle=\bfseries,title={Plain English: #1},boxrule=0.6pt,arc=2pt}
\newtcolorbox{keybox}[1]{breakable,colback=gold,colframe=orange!70!black,fonttitle=\bfseries,title={#1},boxrule=0.6pt,arc=2pt}
\newtcolorbox{warnbox}[1]{breakable,colback=crim!6,colframe=crim,fonttitle=\bfseries,title={#1},boxrule=0.6pt,arc=2pt}

% ---- math shortcuts ------------------------------------------------------
\newcommand{\Ncdf}{\mathrm{N}}          % standard normal CDF
\newcommand{\PD}{\mathrm{PD}}
\newcommand{\DD}{\mathrm{DD}}
\newcommand{\E}{\mathbb{E}}
\newcommand{\Rs}{\textup{Rs}\,}

\lstset{basicstyle=\ttfamily\footnotesize,keywordstyle=\color{navy}\bfseries,commentstyle=\color{teal}\itshape,
  stringstyle=\color{crim},numbers=left,numberstyle=\tiny\color{gray},frame=single,rulecolor=\color{navy!40},
  breaklines=true,showstringspaces=false,language=Python,columns=fullflexible,xleftmargin=1.5em,framexleftmargin=1.5em}

\captionsetup{font=small,labelfont=bf}
\setlist{itemsep=2pt,topsep=4pt}
\setlength{\parskip}{4pt}\setlength{\parindent}{0pt}
\pagestyle{fancy}\fancyhf{}\fancyhead[L]{\small From Committee to Continuum}\fancyhead[R]{\small MS37025 master note}\fancyfoot[C]{\thepage}

\begin{document}

% ---- title block ---------------------------------------------------------
\begin{tcolorbox}[colback=navy,colframe=navy,coltext=white,arc=3pt]
{\Large\bfseries From Committee to Continuum}\\[3pt]
{\large\itshape Can Markets Price Default Faster Than Rating Agencies?}\\[3pt]
Master study note: credit ratings, the Merton/KMV model, and the Yes Bank case
\end{tcolorbox}

\begin{tabularx}{\textwidth}{@{}lX@{}}
\textbf{Name:} & Dhritabrata Swarnakar\\
\textbf{Roll No.:} & 25MSG1R16\\
\textbf{Course:} & MS37025 --- Managing Financial Services and Institutions\\
\textbf{Professor:} & S. L. Tulasi Devi\\
\textbf{Department:} & Department of Management Studies, National Institute of Technology, Warangal\\
\end{tabularx}

\vspace{6pt}
\begin{keybox}{What this project is, in one paragraph}
This is a research project I built end to end: I derived a credit-risk model from first principles, wrote the code that runs it, ran it on Yes Bank's real price data, and then checked what I found against the real public record of what actually happened to Yes Bank between 2018 and 2020. Credit rating agencies (CRISIL, ICRA, CARE) judge a bank's default risk through a committee, at intervals, using what the bank discloses to them. The Merton model does something structurally different: it reads a default signal continuously, every single trading day, straight out of the bank's share price. My question was simple to state and hard to answer honestly: did that market-implied signal, which I call \emph{distance to default}, turn bad before the rating agencies downgraded Yes Bank, and before the RBI stepped in? I answer that with my own computed numbers in Section~\ref{sec:verdict}, I show you exactly how I built the pipeline that produced those numbers in Part~\ref{part:build}, and I spend Part~\ref{part:realworld} on the question that actually matters outside a classroom: so what? What does this tell a regulator, an investor, or a risk manager that they didn't already know?
\end{keybox}

\paragraph{How I've organised this, and how to read it.} Part~\ref{part:theory} is the theory, built from nothing — I assume you know algebra and basic finance, and nothing about option pricing. Part~\ref{part:case} is the Yes Bank case: the real timeline, corrected and verified against primary sources, and my results against it. Part~\ref{part:build} is the part most write-ups skip: the actual build, in the order I did it, with the real numbers my code spat out and what each number meant the moment I saw it. Part~\ref{part:realworld} is the real-world interpretation — not "what does the math say" but "what does this mean for how credit risk actually gets managed." Part~\ref{part:close} is the honest accounting: limitations, what I verified and how, and a glossary. Boxes marked \emph{Plain English} restate a formula in everyday words with a number plugged in; boxes marked \emph{FRM link} note where this shows up in FRM-style study material.

\begin{table}[H]\centering\small
\begin{tabular}{@{}cp{10.8cm}@{}}\toprule
\textbf{Symbol} & \textbf{Meaning (used throughout)}\\\midrule
$V$ & value of the firm's assets today (not directly observable)\\
$D$ & default point: the debt level the assets must cover\\
$E$ & value of the firm's equity (share price $\times$ number of shares)\\
$\mu$ & expected asset drift: the average yearly growth rate of $V$ in the real world\\
$\sigma$ (or $\sigma_V$) & asset volatility: yearly standard deviation of asset returns\\
$\sigma_E$ & equity volatility: yearly standard deviation of equity returns\\
$T$ & time horizon in years (usually 1)\\
$r$ & risk-free interest rate (continuously compounded)\\
$\Ncdf(\cdot)$ & standard normal cumulative distribution function\\
$\DD$ & distance to default, in standard-deviation units\\\bottomrule
\end{tabular}
\end{table}

\tableofcontents
\clearpage

\part{The Theory: Building Merton From Nothing}
\label{part:theory}
\section{The setup: how India assesses bank credit risk today}

Before I touch a single formula, I want to lay out the system I'm arguing against — or rather, the system I'm testing an alternative to. I'm not claiming credit rating agencies are incompetent. I'm claiming their process has a built-in lag, and I wanted to measure that lag, in days, on a real bank, rather than just asserting it.

\subsection{What a credit rating is}
\begin{defbox}{credit rating}
A credit rating is an agency's \emph{opinion}, written as a letter grade, of how likely a borrower is to fail to pay what it owes. In India the main agencies are CRISIL, ICRA and CARE (India Ratings and Brickwork also rate banks). ``Default'' here means missing a payment on a rated obligation. Grades run from AAA (very low default risk) down through the AA, A and BBB bands to D (already in default). BBB$-$ and above is called \emph{investment grade}; below that is \emph{sub-investment grade}, informally ``junk''.
\end{defbox}

A bank has many obligations, so it has many ratings: one for senior bonds, one for subordinated (Tier~II) bonds, another for additional Tier~I (AT-1) bonds, and a short-term rating for certificates of deposit. The ratings differ because the claims differ in how much loss they absorb. AT-1 bonds, for example, can be written down when the bank's capital falls below a trigger, so they are rated several \emph{notches} (one step on the letter scale, e.g.\ AA to AA$-$) below senior debt.

\subsection{How the process works}
\begin{enumerate}
\item \textbf{Evidence gathering.} Analysts read audited financial statements, regulatory filings and management presentations, and meet management. For a bank they focus on capital adequacy (the buffer of own funds against risky assets), asset quality (how many loans are bad), profitability and funding.
\item \textbf{Committee decision.} A rating committee votes on the grade. The output is a judgement, not the result of a formula.
\item \textbf{Surveillance.} Once assigned, the rating is monitored. It is formally reviewed at intervals (at least once a year) and reviewed earlier after significant events. Agencies also signal likely changes through an \emph{outlook} (positive, stable, negative: the direction a rating may move over the medium term) or a \emph{watch} (a shorter-term flag that a decision is pending after a specific event).
\end{enumerate}

\begin{keybox}{Three built-in properties}
\textbf{Committee-based} (judgement, not a formula), \textbf{periodic} (the published grade changes at discrete moments), and \textbf{reliant on disclosure} (agencies see what the bank reports, when it reports it). Each property is a possible source of delay. The next section asks whether the delay is a flaw or a design choice.
\end{keybox}

\begin{defbox}{the school report card vs. the fitness tracker}
Here is the analogy I kept coming back to while building this project. A credit rating is like a school report card: it is issued at fixed points in the term, it reflects a teacher's holistic judgement of the whole term so far, and it deliberately does not change every time a student has one bad week — a single poor quiz shouldn't crash a semester grade that otherwise reflects consistent effort. A fitness tracker, on the other hand, updates every second from a pulse sensor on your wrist. It has no judgement in it at all — it just reports a number, continuously, whether that number means anything in isolation or not. The Merton model I build in this project is the fitness tracker. The rating agencies are the report card. Neither is wrong to exist; they are built to answer different questions, and the whole point of my project is to see how differently they actually behaved on one real patient: Yes Bank.
\end{defbox}

\section{The critique: why ratings are structurally slow}
\label{sec:critique}

\subsection{Ratings are sticky by design}
Agencies describe their ratings as \emph{through-the-cycle}: an assessment meant to hold across the ups and downs of the economy rather than reflect every swing in prices. That is deliberate. Investors and regulators use ratings in mandates and covenants (for example, a fund that may only hold investment-grade bonds), so a rating that flipped with each market move would force repeated buying and selling. Agencies therefore prefer to wait for documented evidence before acting. The result is \emph{stickiness}: the rating moves late and in steps.

\begin{defbox}{sticky rating}
A sticky rating is one that changes slowly and in discrete steps even while the underlying credit quality is changing continuously. Stickiness is not a bug in the agency's process. It follows from a committee that needs evidence, a disclosure calendar that delivers evidence quarterly, and a rating that has real consequences for holders.
\end{defbox}

\subsection{A recurring pattern}
The criticism that agencies react late is not specific to India. Two well-known global episodes are usually cited.
\begin{itemize}
\item \textbf{Enron (2001).} Enron was still rated investment grade in late November 2001. The agencies cut it to sub-investment grade on 28 November, days before its bankruptcy filing on 2 December 2001.
\item \textbf{The 2007--2009 crisis.} Large volumes of mortgage-linked securities carried AAA ratings and were then downgraded in bulk from mid-2007 onward, after mortgage losses had already begun to show.
\end{itemize}
Both are widely documented; they are background here, and the Yes Bank case in Section~\ref{sec:case} is the local test.

\subsection{The Yes Bank illustration}
CARE upgraded Yes Bank's senior and Tier~II debt to AAA on 5 July 2018 \cite{bsCareWatch}. On 21 September 2018 the RBI declined to give the founder-CEO a fresh three-year term \cite{bsIcraWatch}, and the share price fell 29\% that day (in our price data). CARE placed its AAA on credit watch about ten days later, ICRA placed its rating on watch on 17 November \cite{bsIcraWatch}, and the first actual downgrades by ICRA and CARE came on 28 November 2018 \cite{bsNov28}. Between the AAA upgrade date and the day before the first downgrade, the share price had fallen 48\%. That gap between a fast-moving price and a slow-moving rating is the empirical question of this note.

\section{The alternative: Merton's model, built from first principles}
\label{sec:merton}

\subsection{The problem, and why option pricing applies}
Here is what I was actually trying to do: get the probability that a firm cannot pay its debts, using information that updates every single day instead of every few months. The firm's equity trades on an exchange, so I can observe its value and how much it bounces around, every trading day, for free, from any stock exchange's closing price. Merton (1974) \cite{merton1974} made the connection I needed: a firm's equity has the exact same shape of payoff as a financial \emph{option}, so the maths built for pricing options can be turned around and used to read default risk straight out of equity prices. This single observation is the hinge the entire rest of this project swings on, so I want to build it slowly and make sure it actually lands before I write a single line of code.

\begin{defbox}{structural model}
A \emph{structural} credit model explains default from the firm's own balance-sheet structure: the firm defaults when its assets are worth too little to cover its debts. It contrasts with \emph{reduced-form} models, which treat default as a random event with a probability that depends on outside variables and never look at assets versus debt.
\end{defbox}

Set up the simplest firm. It is financed by shareholders' equity and by a single zero-coupon bond: a bond that pays one lump sum, the face value $D$, at a fixed date $T$ and nothing before. At $T$ two things can happen.
\begin{itemize}
\item If the assets are worth more than $D$ ($V_T>D$), the firm repays the debt in full and shareholders keep the remainder $V_T-D$.
\item If the assets are worth less ($V_T<D$), the firm cannot pay. Lenders take over the assets, and shareholders receive nothing (limited liability: they cannot be asked for more).
\end{itemize}
So shareholders receive $\max(V_T-D,\,0)$, and lenders receive $\min(V_T,\,D)$.

\begin{figure}[H]\centering
\includegraphics[width=0.78\textwidth]{figs/payoff.png}
\caption{Payoffs at maturity. Equity is a call option on the firm's assets with strike $D$; debt is capped at $D$ and loses value whenever assets fall short.}
\end{figure}

A \emph{call option} on an asset with \emph{strike price} $K$ pays $\max(S_T-K,0)$ at expiry, where $S_T$ is the asset's price then. Equity has exactly this payoff with $S_T=V_T$ and $K=D$. \textbf{Equity is a call option on the firm's assets, with the debt as the strike price.} Default is the event that the option expires worthless, $V_T<D$. Equivalently, lenders own risk-free debt but have sold shareholders a put option on the assets, and the value of that put is what a credit spread compensates for.

\begin{defbox}{the mortgaged-flat analogy}
I find this easiest to picture with a home loan. Say you buy a flat worth 1 crore using a 70 lakh bank loan, and the loan is structured so the whole amount is due at the end of the term in one shot. You are, economically, in exactly the shareholders' position: you own whatever is left after the loan is repaid. If the flat's market value rises to 1.3 crore by the time the loan is due, you pay off the 70 lakh and keep 60 lakh of profit. But if a flood damages the building and its value crashes to 50 lakh — below what you owe — you are not on the hook for the shortfall. You simply hand the keys to the bank (this is sometimes called ``jingle mail'' in mortgage markets) and walk away owing nothing more, because the loan is secured only against the flat. Your payoff was $\max(\text{flat value} - \text{loan}, 0)$ the whole time, whether you realised it or not — and that is a call option. The bank, meanwhile, was never going to get more than what it was owed even if the flat tripled in value, and it bears the loss if the flat's value falls below the loan — exactly the lender's position in Merton's model.
\end{defbox}

\begin{keybox}{Why this helps}
If we can describe how $V$ moves, then the probability of default is simply the probability that a random variable ends below $D$. The next subsection builds the tools to describe that movement.
\end{keybox}

\subsection{Prerequisites, from scratch}

\subsubsection{The normal distribution and $\Ncdf(\cdot)$}
A random variable $Z$ is \emph{standard normal} if it follows the bell curve with mean~0 and standard deviation~1. The function $\Ncdf(x)=\Pr(Z\le x)$ gives the probability that $Z$ falls at or below $x$. Useful anchors: $\Ncdf(0)=0.5$, $\Ncdf(-1)\approx0.159$, $\Ncdf(-2)\approx0.023$, $\Ncdf(-3)\approx0.0013$. Going one standard deviation into the left tail leaves a 16\% chance; three leaves about 0.13\%. This mapping from ``number of standard deviations'' to ``probability'' is the engine behind distance to default.

\subsubsection{Why asset values are lognormal, not normal}
Suppose asset value could be modelled as a normal variable. It could then go negative, but a firm's assets cannot be worth less than zero, and shocks act in \emph{percentage} terms (a 10\% fall hurts a large firm more in rupees than a small one). So we model the \emph{log} of asset value as normal. A variable whose logarithm is normal is called \emph{lognormal}. It stays positive, and its distribution is skewed: a long tail to the upside and a bounded floor at zero.

\subsubsection{Geometric Brownian motion (GBM)}
GBM is the standard model for a quantity that grows by random percentage steps. Over a tiny time step $dt$,
\begin{equation}
\frac{dV_t}{V_t} \;=\; \mu\,dt \;+\; \sigma\,dW_t ,
\label{eq:gbm}
\end{equation}
where
\begin{itemize}
\item $dV_t/V_t$ is the \emph{percentage change} in asset value during the step;
\item $\mu$ is the \textbf{drift}: the average yearly percentage growth. If $\mu=8\%$, assets grow on average 8\% a year, ignoring randomness;
\item $\sigma$ is the \textbf{volatility}: the standard deviation of the yearly percentage change. If $\sigma=25\%$, a typical yearly swing around the average is about 25 percentage points;
\item $dW_t$ is a \emph{Wiener increment}: a random shock, normally distributed with mean~0 and variance $dt$, independent from step to step.
\end{itemize}
Solving equation~\eqref{eq:gbm} (this uses It\^o's lemma, the calculus rule for random processes) gives the value at time $T$:
\begin{equation}
V_T \;=\; V_0\,\exp\!\Big[\big(\mu-\tfrac12\sigma^2\big)T+\sigma\sqrt{T}\,Z\Big],\qquad Z\sim\text{N}(0,1).
\label{eq:gbmsol}
\end{equation}
Taking logs, $\ln V_T$ is normal with mean $\ln V_0+(\mu-\tfrac12\sigma^2)T$ and standard deviation $\sigma\sqrt T$. That is exactly the lognormal shape described above.

\begin{warnbox}{The $\mu-\tfrac12\sigma^2$ term is not CAGR}
Equation~\eqref{eq:gbmsol} contains $\mu-\tfrac12\sigma^2$, not $\mu$. The reason is \emph{volatility drag}. Take an asset worth 100 that rises 50\% and then falls 50\%. The arithmetic average of the two returns is $0\%$, yet the asset is worth $100\times1.5\times0.5=75$. Randomness costs growth, and the cost is about $\tfrac12\sigma^2$ per year.

Precisely: the \emph{mean} of $V_T$ grows at $\mu$, since $\E[V_T]=V_0e^{\mu T}$, but the \emph{median} (the typical outcome) grows at $\mu-\tfrac12\sigma^2$, since the median of $V_T$ is $V_0e^{(\mu-\frac12\sigma^2)T}$. The term is the drag already built into the formula, subtracted once.

\textbf{Why substituting CAGR for $\mu$ is wrong.} The compound annual growth rate (CAGR) is the growth rate that turns the start value into the realised end value. It is a geometric average, and it \emph{already contains} the volatility drag: over long samples, CAGR $\approx\mu-\tfrac12\sigma^2$. Replace $\mu$ by CAGR in the formula and the drag is subtracted twice: once when the CAGR was measured, and again by $-\tfrac12\sigma^2$. The input must be the arithmetic expected return $\mu$.
\end{warnbox}

For our illustrative firm ($V_0=120$, $\mu=8\%$, $\sigma=25\%$, $T=1$) the mean end value is $120e^{0.08}=130.0$, while the median is $120e^{0.08-0.03125}=126.0$.

\subsubsection{Where $\Ncdf(d_2)$ comes from in Black--Scholes}
The Black--Scholes formula (1973) \cite{blackScholes1973} prices a European call option (exercisable only at expiry). Its key idea is to price by \emph{no-arbitrage} (no risk-free profit): one can build a portfolio of the asset and cash that copies the option's payoff, so the option's price equals the cost of that portfolio. This pricing does not depend on investors' appetite for risk, so it can be computed in an imagined \emph{risk-neutral world}, where every asset grows at the risk-free rate $r$ instead of its real-world drift.

\begin{defbox}{risk-neutral}
``Risk-neutral'' means priced as if investors were indifferent to risk and asked for no extra return for bearing it. In that imagined world all assets grow at $r$. It is a computational device that gives the right \emph{price}; it does not describe how assets really grow. Probabilities computed in it are called risk-neutral probabilities.

Here's the version of this I actually hold onto: imagine two insurance agents valuing the same fire-insurance policy on a warehouse. One agent genuinely believes the warehouse has a 2\% chance of burning down this year — that's the real-world, physical probability. The other agent doesn't care what the true odds are; she just wants to quote a price that makes the insurer indifferent between selling the policy and not selling it, given what she can hedge in the market. Her quoted price might imply a 3\% chance of fire, not because she thinks fires are more likely, but because she's baked in a margin for the risk she's taking on. That 3\% is her risk-neutral number — useful for pricing, misleading if you mistake it for her actual belief about fire risk.
\end{defbox}

For a call with strike $K$ on an asset worth $S$, the Black--Scholes value is $S\,\Ncdf(d_1)-Ke^{-rT}\Ncdf(d_2)$, where
\[
d_1=\frac{\ln(S/K)+\big(r+\tfrac12\sigma^2\big)T}{\sigma\sqrt T},\qquad d_2=d_1-\sigma\sqrt T .
\]
The term $\Ncdf(d_2)$ has a clean meaning: it is the \textbf{risk-neutral probability that the option finishes in the money}, i.e.\ that $S_T>K$. The second term of the price, $Ke^{-rT}\Ncdf(d_2)$, is then the present value of paying $K$ times the probability that you have to. The term $\Ncdf(d_1)$ is the option's \emph{delta}: how many rupees the option value moves when the asset moves by one rupee. We will use both.

\subsection{Merton's model: equity and debt values}
Assume $V$ follows GBM, the firm has one zero-coupon bond with face $D$ due at $T$, markets are frictionless, and default can happen only at $T$. Equity is a call option on $V$ with strike $D$, so by Black--Scholes
\begin{equation}
E \;=\; V\,\Ncdf(d_1)\;-\;D\,e^{-rT}\,\Ncdf(d_2),\qquad
d_1=\frac{\ln(V/D)+\big(r+\tfrac12\sigma_V^2\big)T}{\sigma_V\sqrt T},\quad d_2=d_1-\sigma_V\sqrt T .
\label{eq:equity}
\end{equation}
Here $V,D,T,r$ are as in the symbol table and $\sigma_V$ is asset volatility. Debt is whatever assets are not equity: $B=V-E=V\,\Ncdf(-d_1)+D\,e^{-rT}\Ncdf(d_2)$. The market's implied yield on this debt is $y=-\tfrac1T\ln(B/D)$, and the \emph{credit spread} (the extra yield over the risk-free rate) is
\begin{equation}
y-r \;=\; -\frac1T\ln\!\Big[\Ncdf(d_2)+\frac{V}{De^{-rT}}\Ncdf(-d_1)\Big].
\end{equation}

\paragraph{Worked example.} A firm has $V=120$, $D=100$, $T=1$, $\sigma_V=25\%$, $r=6\%$. Then $d_2=0.844$, $d_1=1.094$, $\Ncdf(d_1)=0.863$, $\Ncdf(d_2)=0.801$, so
\[
E=120(0.863)-100e^{-0.06}(0.801)=103.57-75.41=28.16,
\]
debt is worth $B=120-28.16=91.84$, the implied yield is $8.51\%$ and the credit spread is $2.51\%$. Leverage magnifies volatility: asset volatility of 25\% becomes equity volatility of $\Ncdf(d_1)\sigma_VV/E=92\%$.

\subsection{The distance-to-default formula}
\label{sec:dd}
The risk-neutral world prices debt. To describe the chance of default in the \emph{real} world, we return to real-world drift $\mu$. Default at $T$ means $V_T<D$, i.e.\ $\ln V_T<\ln D$. Since $\ln V_T$ is normal with the mean and standard deviation found above,
\[
\Pr(V_T<D)=\Ncdf\!\left(\frac{\ln D-\ln V_0-(\mu-\tfrac12\sigma^2)T}{\sigma\sqrt T}\right)=\Ncdf(-\DD),
\]
where
\begin{equation}
\boxed{\;\DD \;=\; \frac{\ln(V/D)+\big(\mu-\tfrac12\sigma^2\big)T}{\sigma\sqrt T}\;}
\qquad\text{and}\qquad \PD=\Ncdf(-\DD).
\label{eq:dd}
\end{equation}

\begin{table}[H]\centering\small
\begin{tabular}{@{}cp{11cm}@{}}\toprule
\textbf{Symbol} & \textbf{What it is, in plain words}\\\midrule
$V$ & Market value of the firm's assets today. Not observed; estimated (Section~\ref{sec:estimate}).\\
$D$ & Default point: the debt level the assets must cover. For a bank this choice is contested (Section~\ref{sec:limits}).\\
$\ln(V/D)$ & Today's cushion, in logs: how far the assets stand above the default point in percentage terms. Zero means assets equal debt.\\
$\mu$ & Expected asset drift: average yearly growth of $V$ (arithmetic mean, not CAGR).\\
$\tfrac12\sigma^2$ & Volatility drag: the compounding penalty that randomness imposes.\\
$T$ & Horizon in years.\\
$\sigma$ & Asset volatility: yearly standard deviation of asset returns.\\
$\sigma\sqrt T$ & The size of one standard deviation of the log-asset value over the horizon (uncertainty grows with the \emph{square root} of time).\\\bottomrule
\end{tabular}
\end{table}

\begin{defbox}{distance to default}
The numerator is where we expect the log-asset value to sit above the default point at time $T$: today's cushion plus the expected drift. The denominator is the size of one standard deviation of uncertainty. So $\DD$ counts how many standard deviations separate the expected asset value from the default point. A $\DD$ of 4 means default needs a four-standard-deviation drop (rare); a $\DD$ near 0 means the expected outcome is right at the default point, so default is about a coin-flip.

The analogy I use: think of $\DD$ as a car's following distance on a highway, measured not in metres but in ``how many of my own braking distances am I currently keeping.'' A following distance of 4 braking-distances is comfortable — you would need a very sudden, very hard stop from the car ahead to hit it. A following distance of 0.5 braking-distances means you are basically tailgating — almost any unexpected stop ends in a collision. $\DD$ is the same idea applied to a company's balance sheet: it is not a distance in rupees, it is a distance measured in units of ``how much random wobble this specific firm's assets typically have.'' That is exactly why I cannot compare a $\DD$ of 2 for a volatile startup with a $\DD$ of 2 for a stable public-sector bank and call them equally safe — the units are each firm's own volatility, not a common ruler.
\end{defbox}

\paragraph{Worked example (continued).} With $V=120$, $D=100$, $\mu=8\%$, $\sigma=25\%$, $T=1$:
\[
\DD=\frac{\ln1.2+(0.08-0.03125)}{0.25}=\frac{0.1823+0.0488}{0.25}=0.924,\qquad \PD=\Ncdf(-0.924)=17.8\%.
\]
Section~\ref{sec:mc} confirms this by simulation.

\subsection{Physical PD versus risk-neutral PD}
\label{sec:physical-rn}
Two default probabilities come out of the same model, and they differ.
\begin{align}
\PD_{\text{physical}}&=\Ncdf(-\DD),\quad \DD\text{ uses the real-world drift }\mu,\\
\PD_{\text{risk-neutral}}&=\Ncdf(-d_2),\quad d_2=\frac{\ln(V/D)+\big(r-\tfrac12\sigma^2\big)T}{\sigma\sqrt T}\ \ (\text{same as }\DD\text{ with }\mu\to r).
\end{align}
The only difference is the drift: real assets are expected to grow at $\mu$, but in the risk-neutral world they grow at $r$. Investors demand a premium ($\mu>r$) for holding risky assets, so real-world growth is higher and real-world default is less likely: $\PD_{\text{physical}}<\PD_{\text{risk-neutral}}$. In the worked example, $d_2=0.844$ and the risk-neutral PD is 19.9\% against the physical 17.8\%. Use the risk-neutral PD for \emph{pricing} bonds and credit derivatives; use the physical PD to estimate how likely default \emph{actually is}.

\subsection{Estimating $V$ and $\sigma_V$ in practice}
\label{sec:estimate}
Asset value $V$ and asset volatility $\sigma_V$ are not quoted anywhere: assets (loans, securities, buildings) do not trade as a single price. What we can observe is equity value $E$ (share price times shares outstanding) and equity volatility $\sigma_E$ (the standard deviation of daily log-returns, scaled by $\sqrt{252}$ to a yearly figure). The model links the two sets of quantities by two equations.

\begin{keybox}{Two equations, two unknowns}
\begin{align}
E &= V\,\Ncdf(d_1)-D\,e^{-rT}\,\Ncdf(d_2) &&\text{(equity is a call on assets)}\label{eq:e1}\\
\sigma_E\,E &= \Ncdf(d_1)\,\sigma_V\,V &&\text{(equity volatility is asset volatility, scaled by leverage)}\label{eq:e2}
\end{align}
Given observed $E$, $\sigma_E$ and the balance-sheet input $D$, solve numerically for $V$ and $\sigma_V$.
\end{keybox}
Equation~\eqref{eq:e2} follows because a small change in $V$ changes $E$ by $\Ncdf(d_1)$ times as much (the delta), so the percentage volatility of $E$ equals the volatility of $V$ scaled by $\Ncdf(d_1)V/E$. Solving is a root-finding problem: pick a trial $\sigma_V$, solve~\eqref{eq:e1} for $V$, check whether~\eqref{eq:e2} holds, and adjust $\sigma_V$ until it does. Our code does this with two nested root finders. In the worked example, feeding in $E=28.16$ and $\sigma_E=92\%$ returns $V=120$ and $\sigma_V=25\%$.

\begin{warnbox}{The model does not predict default ``exactly''}
Asset value and asset volatility are estimated, not observed. $E$ and $\sigma_E$ carry market noise; $D$ is a modelling choice; $r$ and $\mu$ are assumptions; and the normal-distribution shape is itself an assumption. Every $\DD$ in this note is an estimate with error around it. It is a risk gauge, not a forecast of a date.
\end{warnbox}

\subsection{KMV's practical refinement}
Moody's KMV (the commercial descendant of Merton's model, built by Kealhofer, McQuown and Vasicek) \cite{crosbie2003} made two changes for practice.
\begin{enumerate}
\item \textbf{The default-point convention.} Real firms do not default the moment assets dip below total liabilities; they often keep operating while long-dated debt is not yet due. KMV set the default point at \emph{short-term debt plus one half of long-term debt} for a one-year horizon. This is an empirical convention, not a theorem.
\item \textbf{DD $\to$ EDF.} The normal distribution understates how often firms actually default (real asset returns have fat tails and jumps). KMV therefore does not use $\Ncdf(-\DD)$. It ranks firms by $\DD$ and records, from its history of real defaults, what fraction of firms with a given $\DD$ defaulted within a year. That observed fraction is the \emph{expected default frequency} (EDF). $\DD$ is the ranking; EDF is the calibrated probability.
\end{enumerate}

\begin{keybox}{FRM link}
Structural models (Merton, KMV) appear in FRM credit-risk study as the bridge from balance-sheet leverage to default probability; the equity-as-call insight and the DD to EDF mapping are the parts most likely to be examined conceptually. See Hull and Malz \cite{hull,malz2011}.
\end{keybox}


\part{The Case: Yes Bank, 2018--2020}
\label{part:case}
\section{The mechanics: a Monte Carlo simulation of the model}
\label{sec:mc}

Before I trusted my own derivation enough to run it on a real bank, I wanted to prove it to myself the hard way. The closed-form $\PD=\Ncdf(-\DD)$ is a statement about a random variable. A \emph{Monte Carlo simulation} tests such a statement by brute force: generate many random futures, count how often the event occurs, and see whether the count matches the formula. This section is me checking my own homework before I let the model anywhere near Yes Bank's actual share price.

\begin{defbox}{Monte Carlo simulation}
Draw a large number of random scenarios from the model, record for each whether default occurred, and take the fraction of scenarios with default as the estimated probability. By the law of large numbers the estimate converges to the true probability, with a standard error of about $\sqrt{p(1-p)/n}$ for $n$ scenarios.
\end{defbox}

\paragraph{Algorithm.} (1) Fix $V_0$, $D$, $\mu$, $\sigma$, $T$. (2) For each scenario draw $Z\sim\text{N}(0,1)$ and compute $V_T$ from equation~\eqref{eq:gbmsol}. (3) Count the fraction of scenarios with $V_T<D$. (4) Compare with $\Ncdf(-\DD)$. A second version simulates the whole path in daily steps and counts a default if the path ever touches $D$ (a \emph{barrier} or first-passage variant, following Black and Cox \cite{blackcox1976}). That variant gives a higher probability because a path can dip below $D$ and recover before $T$; Merton's original model counts default only at $T$.

\begin{figure}[H]\centering
\includegraphics[width=0.62\textwidth]{figs/gbm_paths.png}\hfill
\includegraphics[width=0.36\textwidth]{mc_convergence.png}
\caption{Left: 60 simulated GBM asset paths for the illustrative firm; red paths end below the default point $D=100$. Right: the Monte Carlo default probability converges to the closed-form value as the number of paths grows.}
\end{figure}

\lstinputlisting[caption={Monte Carlo check of the Merton default probability (file \texttt{merton\_monte\_carlo.py}). Run in VS Code after \texttt{pip install numpy scipy matplotlib}.},label={lst:mc}]{merton_monte_carlo.py}

\paragraph{Output of the listing.}
\begin{center}\small
\begin{tabular}{rrrr}\toprule
\textbf{Paths} & \textbf{MC PD} & \textbf{95\% CI $\pm$} & \textbf{Abs.\ error}\\\midrule
100 & 13.00\% & 6.59\% & 4.77\%\\
1,000 & 18.80\% & 2.42\% & 1.03\%\\
10,000 & 18.10\% & 0.75\% & 0.33\%\\
100,000 & 17.83\% & 0.24\% & 0.06\%\\
1,000,000 & 17.71\% & 0.07\% & 0.06\%\\\midrule
\multicolumn{4}{c}{Closed form $\Ncdf(-\DD)$ with $\DD=0.9243$: \textbf{17.77\%}. Barrier variant (100k paths): 37.7\%.}\\\bottomrule
\end{tabular}
\end{center}
The simulated PD converges to the formula, and the errors sit inside the confidence interval. (Exact figures for small $n$ depend on the random seed.)

\section{The case: Yes Bank, 2018--2020}
\label{sec:case}

This is the part of the project where I stopped working with an illustrative firm and started working with a real one. Everything from here on uses Yes Bank's actual traded price, Yes Bank's actual published balance sheets, and the actual dated public record of what the rating agencies and the RBI did. I went back after my first draft and verified every date I was unsure of against a primary source — a regulatory filing, a press release, or a SEBI order — rather than leaving them as things I could not check. I'll flag the handful of places where I still could not get a primary source, but the list is much shorter now than when I started.

\subsection{Timeline from the public record, verified}
Sources are numbered in the bibliography. A date I verified against a company filing, an agency's own press release, or a regulator's order is given without a qualifier. A date that still rests only on secondary news reporting is marked so, honestly, rather than dressed up as something I confirmed.

\begingroup\small
\renewcommand{\arraystretch}{1.15}
\begin{longtable}{@{}>{\raggedright\arraybackslash}p{2.5cm}p{11.1cm}@{}}
\toprule\textbf{Date} & \textbf{Event}\\\midrule\endhead
12 May 2017 & Annual report discloses RBI-assessed FY16 gross NPAs of Rs 4,925.68 crore against Rs 748.98 crore reported: a divergence of Rs 4,176.70 crore \cite{bsDivFY16}.\\
5 Jul 2018 & CARE upgrades senior and Tier II debt to AAA from AA+ --- the first AAA any agency had given Yes Bank, and only the fourth private bank in India (after HDFC, ICICI and Kotak) to hold it \cite{bsCareWatch,sebiCareOrder}. Share price: Rs 349.\\
20 Aug 2018 & Share price peaks (closing high Rs 394 in our data; intraday high Rs 404).\\
17 Sep 2018 & \textbf{Verified.} The RBI's letter curtailing Rana Kapoor's term as MD \& CEO to 31 January 2019 is dated this day. The bank received it the same day, but did not disclose it to the exchanges until the evening of 19 September \cite{bsSep19disc}.\\
19--21 Sep 2018 & The disclosure becomes public on 19 September (a Wednesday); the stock falls across the next two trading sessions, with the sharpest single-day fall (around 29--34\%, depending on the exchange and the moment measured) on 21 September \cite{bsSep21fall}. My own price series has no entry for 20 September.\\
1 Oct 2018 & \textbf{Verified.} CARE places the AAA rating on credit watch with developing implications, citing the RBI's communication on Kapoor's tenure \cite{bsCareOct1}.\\
17 Nov 2018 & ICRA places Yes Bank on watch with negative implications \cite{bsIcraWatch}.\\
27 Nov 2018 & Moody's cuts the foreign-currency issuer rating to Ba1 from Baa3 \cite{bsNov28}.\\
28 Nov 2018 & ICRA downgrades senior debt to AA (from AA+) and subordinated to AA$-$ (from AA); CARE downgrades senior to AA+ (from AAA) and subordinated to AA (from AA+) \cite{bsNov28}.\\
24 Jan 2019 & Q3 FY19 results; RBI approves Ravneet Gill as MD \& CEO, to join on or before 1 Mar 2019 \cite{bsQ3FY19}.\\
13 Feb 2019 & Bank states that the RBI found no divergence for FY18; the same report recalls the FY17 divergence (RBI-assessed gross NPAs Rs 8,373.8 crore against Rs 2,018 crore reported) \cite{bsDivFY18}.\\
30 Apr 2019 & Share falls about 29\% on the first session after the first-ever quarterly loss (Rs 1,506 crore in Q4 FY19) \cite{bsFall19}.\\
3 May 2019 & ICRA downgrades Tier I bonds to A (from AA$-$) and Tier II to AA$-$ (from AA), outlook negative \cite{bsIcraMay19,icraJul19}.\\
8 May 2019 & \textbf{Verified.} India Ratings cuts the long-term issuer rating to IND~AA$-$ from IND~AA+, negative outlook; filed late evening, reported the next morning \cite{bsIndraMay19}.\\
24 Jul 2019 & ICRA downgrades a Rs 32,911.7 crore bond programme (one notch on Rs 22,111.7 crore; two notches on Rs 10,800 crore of AT-1) \cite{icraJul19,bsIcraJul19}.\\
15 Aug 2019 (approx.) & QIP (a share sale to institutional investors) raises Rs 1,930 crore \cite{ycQ2}.\\
31 Aug 2019 & India Ratings downgrades further, to IND~A+ from IND~AA$-$, negative outlook --- citing slower-than-expected resolution of stressed assets and a smaller-than-hoped capital raise even after the QIP \cite{bsIndraAug19}. I had not found this rating action when I first wrote this note; I only caught it on re-verification.\\
1 Nov 2019 & Q2 FY20 results: gross NPA 7.39\%; deposits Rs 2,09,497 crore; a US\$1.2 bn investment offer noted \cite{ycQ2}.\\
13--14 Nov 2019 & CARE cuts infrastructure and Tier II bonds to A+ (from AA$-$) and AT-1 to BBB+ (from A$-$), on credit watch \cite{bsCareNov19}.\\
19 Nov 2019 & Bank discloses an RBI-assessed FY19 divergence of Rs 3,277 crore in gross NPAs \cite{bsDiv19}.\\
19 Dec 2019 & India Ratings cuts further to IND~A from IND~A+, placed on rating watch negative \cite{bsIndraDec19}.\\
31 Dec 2019 & CARE cuts infrastructure bonds to A, outlook negative \cite{bsCareDec19}.\\
12 Feb 2020 & \textbf{Verified, newly added.} India Ratings cuts a third time, to IND~A$-$ from IND~A, still on rating watch negative --- citing the continued delay and inconclusive size of the hoped-for equity infusion. This is Ind-Ra's third downgrade of Yes Bank inside twelve months \cite{bsIndraFeb20}.\\
5 Mar 2020 & RBI moratorium: withdrawals capped at Rs 50,000, board superseded \cite{btQ3}.\\
6 Mar 2020 & ICRA rates debt worth Rs 52,611.7 crore [ICRA]~D \cite{bsMar9,icraMar20}.\\
13 Mar 2020 & Reconstruction Scheme, 2020 notified under Section 45 of the Banking Regulation Act, 1949; Q3 FY20 results published, showing a quarterly loss of about Rs 18,560 crore \cite{btQ3,crisilMar20,yesReconScheme}.\\
18 Mar 2020 & Moratorium lifted \cite{crisilMar20}. SBI's board approves an eventual total investment of Rs 7,250 crore (an initial tranche of Rs 6,050 crore gave it about 48.2\%); the Scheme capped SBI's stake at 49\% and required it to hold at least 26\% for three years from infusion --- a lock-in that, as it happens, ran until 13 March 2023 \cite{sbiStakeScheme,yesReconScheme}.\\
19 Mar 2020 & CRISIL assigns CRISIL A2 to the bank's certificates of deposit, citing expected systemic support \cite{crisilMar20}.\\
\bottomrule
\end{longtable}
\endgroup

\paragraph{CRISIL.} I confirmed this properly on re-checking: I found no CRISIL rating action on Yes Bank's debt between 2018 and the moratorium. CRISIL's first visible action is the CD rating on 19 March 2020, after the rescue. The pre-crisis comparison below therefore rests on ICRA, CARE, India Ratings and Moody's — which, as the corrected timeline above shows, between them actually produced \emph{seven} distinct downgrade or watch actions across 2018--2020, not the four or five I'd originally tracked. India Ratings alone downgraded Yes Bank three separate times in a twelve-month span (May 2019, December 2019, February 2020) — a detail I missed in my first pass and only found on verification. That matters for my argument: it means the ``stickiness'' I'm describing is not a case of the agencies doing nothing. They were active. They were just consistently behind the market's own signal, as the results in Section~\ref{sec:verdict} show.

\paragraph{A correction I want to be upfront about: the CARE/SEBI interference finding.} When I first drafted this note, I had been told — and initially wrote down — that SEBI had penalised CARE Ratings specifically for delaying its Yes Bank downgrades. Having now read SEBI's actual 2023 adjudication order in this matter \cite{sebiCareOrder}, I need to correct that, because the real finding is more specific and, frankly, more interesting than the version I started with. SEBI did commission a forensic audit (by EY, submitted February 2020) into whistle-blower allegations that CARE's former MD, Rajesh Mokashi, and former non-executive Chairman, S.B. Mainak, had interfered in CARE's rating process for three stressed borrowers: DHFL, IL\&FS, and Yes Bank. The order that followed, after a multi-year investigation including cross-examination of CARE staff, found \textbf{sufficient evidence to sustain the interference charge for DHFL} — WhatsApp messages showed CARE's own rating team and rating-committee members expressing open frustration that they were not being allowed to downgrade DHFL on the timeline they wanted, and the order describes repeated delays traced to the MD's direct involvement. For \textbf{Yes Bank specifically, the order concluded the opposite}: it states in paragraph 79 that ``sufficient material has not been brought to record to sustain the charge of interference ... in the case of rating assigned to YBL [Yes Bank Ltd.]'' One WhatsApp message from a CARE analyst to another, written in the middle of an unrelated DHFL conversation, does say ``Yes bank me bhi aisa hi pressure tha'' (``there was similar pressure at Yes Bank too'') — but SEBI treated this as hearsay, since the analyst was not a direct party to whatever pressure he was describing, and declined to rely on it. \textbf{So the honest statement is:} CARE's internal conduct around DHFL was found, by a regulator, after due process, to have delayed a downgrade that its own analysts wanted to make sooner. For Yes Bank, the same pattern was alleged, smelled plausible to a whistle-blower, and was investigated at length — but was not proven. I think this distinction matters for a research project, not a rumour mill, and I would rather report the finding SEBI actually reached than the stronger, cleaner story I was originally given.

\subsection{What does ``default'' mean here?}
Yes Bank never missed a senior bond payment before the intervention. The trigger was \emph{regulatory}: a moratorium was imposed on 5 March 2020 under Section~45 of the Banking Regulation Act, 1949 \cite{crisilMar20}. Withdrawals were capped at Rs 50,000, the RBI superseded the board \cite{btQ3}, and a rescue was assembled within days. The Merton model defines default as assets falling below debt; the RBI acted on a supervisory view that the bank could not meet its liabilities as they fell due.

For this test I therefore use \emph{three} events, in order: (i) rating actions by the agencies, (ii) the RBI moratorium of 5 March 2020 as the ``default-equivalent'' event, and (iii) ICRA's [ICRA]~D of 6 March 2020, which is the agencies' own default label and came a day after the RBI moved. The question asked of the distance to default is: \emph{did it deteriorate before (i) and before (ii)?}

\subsection{Data and method}
\paragraph{Prices.} The daily closing prices you supplied (1 Jan 2018 -- 30 Apr 2020, 573 trading days), cut at 6 March 2020, the first session after the moratorium news. Equity value is $E=\text{price}\times\text{shares}$, with 231.0 crore shares until mid-August 2019 and 254.9 crore afterwards (after the QIP). These share counts come from book value per share in the bank's press release \cite{ycQ2}: shareholders' funds of Rs 27,331 crore at Rs 118.4 per share (Sep 2018), Rs 26,495 crore at Rs 114.3 (Jun 2019) and Rs 27,790 crore at Rs 109.0 (Sep 2019).

\paragraph{Default point $D$.} $D$ is total liabilities: deposits, borrowings and other liabilities. I used total liabilities rather than KMV's ``short-term plus half of long-term'' because the disclosure needed to split a bank's liabilities by maturity is not available for every quarter, and because deposits, though legally short-term, are the bank's core funding (the ambiguity is discussed in Section~\ref{sec:limits}). Sensitivity to scaling $D$ down is reported below.

\begin{table}[H]\centering\small
\begin{tabular}{@{}lrrrrl@{}}\toprule
\textbf{Balance sheet} & \textbf{Deposits} & \textbf{Borrowings} & \textbf{Other} & \textbf{$D$ used} & \textbf{Public from}\\\midrule
Dec 2017 & 171,731 & 56,302 & 10,500$^\dagger$ & 238,533 & 1 Jan 2018$^\dagger$\\
Mar 2018 & 200,738 & 74,894 & 11,056 & 286,688 & 1 May 2018$^\dagger$\\
Jun 2018 & 213,395 & 78,790 & 12,500$^\ddagger$ & 304,685 & 31 Jul 2018$^\dagger$\\
Sep 2018 & 222,838 & 101,660 & 19,818 & 344,316 & 30 Oct 2018$^\dagger$\\
Dec 2018 & 222,758 & 107,691 & 18,900$^\ddagger$ & 349,349 & 24 Jan 2019\\
Mar 2019 & 227,558 & 108,424 & 17,990 & 353,972 & 27 Apr 2019\\
Jun 2019 & 225,902 & n/a & n/a & 344,666 & 30 Jul 2019$^\dagger$\\
Sep 2019 & 209,497 & n/a & n/a & 318,786 & 1 Nov 2019\\\bottomrule
\end{tabular}
\caption*{\footnotesize Rs crore. $D$ = total balance sheet less shareholders' funds where the split was not available (Sep 2018, Jun 2019, Sep 2019); Mar 2019 is on a consolidated basis. $^\dagger$~my assumption (about 30 days after quarter end, or an approximation of the ``other liabilities'' line); $^\ddagger$~interpolated. The December 2019 accounts were published only on 13 March 2020, so the model, like the market, never saw them before the moratorium.}
\end{table}
Using the \emph{published-at-the-time} balance sheet, rather than the quarter-end figure, mimics what a market participant could see. Sources: \cite{angel3q19,ycQ2,screenerYes}.

\paragraph{Volatility, rate, drift.} $\sigma_E$ is the standard deviation of the previous 60 trading days' log-returns, times $\sqrt{252}$. The risk-free rate is 6\% and the horizon $T=1$ year. Physical $\DD$ uses an asset drift $\mu=10\%$; $d_2$ uses $\mu=r$. Both are assumptions, so both series are plotted and a robustness table follows. $V$ and $\sigma_V$ are solved from equations~\eqref{eq:e1}--\eqref{eq:e2}. The script is \texttt{yes\_bank\_dd.py}.

\subsection{Results}
\begin{figure}[H]\centering
\includegraphics[width=\textwidth]{figs/dd_timeline.png}
\caption{Distance to default (blue: drift $\mu=10\%$; green: $d_2$, drift $=r$) against rating actions and RBI events. The series starts once 60 days of prices exist.}
\label{fig:dd}
\end{figure}

\begingroup\small
\begin{tabular}{@{}p{3.9cm}p{4.8cm}p{6.2cm}@{}}\toprule
\textbf{Episode} & \textbf{Market signal (DD)} & \textbf{Agency response}\\\midrule
CEO-term ruling, Sep 2018 & DD falls from 4.2 to 1.5 on 21 Sep 2018 (baseline median 4.3) & CARE watch $\approx$1 Oct (+10 days); ICRA watch 17 Nov (+57); first downgrades 28 Nov (+68); DD $\approx0.9$ by then\\
Q4 FY19 loss, Apr 2019 & DD 1.9 $\to$ 0.9 on 30 Apr 2019 & ICRA downgrade 3 May (+3 days); India Ratings $\approx$8 May (+8)\\
Slide into the FY19 divergence and capital delay, 2019 & DD first below 0.5 on 3 Oct 2019; low of 0.15 on 5 Nov & CARE cut 13--14 Nov (+42 days); India Ratings 19 Dec (+77)\\
Moratorium, Mar 2020 & DD 3.45 on 4 Mar 2020: highest since Sep 2018; falls to $-0.57$ only on 6 Mar & RBI acts 5 Mar; ICRA [ICRA] D on 6 Mar\\\bottomrule
\end{tabular}
\endgroup

\paragraph{Reading the figure.} There are three features.
\begin{enumerate}
\item \textbf{Every rating episode was preceded, or accompanied, by a market signal.} DD moved on the day of the news, while ratings moved after a lag of days to about two months.
\item \textbf{DD sat far below its 2018 level for the rest of the period.} It was below 60\% of its Mar--Aug 2018 median on 93\% of trading days between 21 Sep 2018 and 4 Mar 2020. The market treated Yes Bank as distressed for 17 months. That is a chronic reading, not a date.
\item \textbf{DD did not flag the intervention.} In January--March 2020 DD \emph{rose} to about 3.5. Two things drove this, both instructive. First, the October 2019 volatility spike left the 60-day window in early January (marked in Figure~\ref{fig:dd}), so measured $\sigma_E$ fell. Second, with equity now only about 3\% of assets, the drift term $\mu\,T$ dominates the tiny cushion $\ln(V/D)$, making physical DD sensitive to the assumed $\mu$ (the risk-neutral $d_2$ rises far less). Also, the December 2019 accounts, with a loss of about Rs 18,560 crore, were not public.
\end{enumerate}

\paragraph{Robustness.} I re-ran the series changing the inputs (ratio to the DD baseline median; a value of 1 means the same as early 2018).
\begin{center}\footnotesize
\begin{tabular}{@{}p{5.2cm}cccccc@{}}\toprule
\textbf{Variant} & \textbf{First flag$^\S$} & \textbf{26 Apr 19} & \textbf{30 Apr 19} & \textbf{31 Oct 19} & \textbf{4 Mar 20}\\\midrule
Base ($r=6\%$, $\mu=10\%$, 60-day, $D\times1$) & 21 Sep 18 & 0.43 & 0.22 & 0.04 & 0.80\\
$\mu=r$ ($d_2$) & 21 Sep 18 & 0.39 & 0.17 & $-0.03$ & 0.34\\
$\mu=15\%$ & 21 Sep 18 & 0.47 & 0.26 & 0.10 & 1.22\\
30-day window & 21 Sep 18 & 0.72 & 0.15 & $-0.09$ & 0.80\\
120-day window & 21 Sep 18 & 0.44 & 0.32 & 0.17 & 0.36\\
$D\times0.85$ & 21 Sep 18 & 0.42 & 0.21 & 0.03 & 0.72\\
$D\times0.70$ & 21 Sep 18 & 0.41 & 0.20 & 0.02 & 0.64\\
$r=5\%$ / $r=7\%$ & 21 Sep 18 & 0.44 / 0.42 & 0.23 / 0.20 & 0.05 / 0.02 & 0.90 / 0.69\\\bottomrule
\end{tabular}

{\footnotesize $^\S$ First date on which DD falls below 60\% of its own Mar--Aug 2018 median.}
\end{center}
The first warning date, the April 2019 drop and the October 2019 trough are the same in all variants. In every variant DD on 4 March 2020 stays above its October 2019 level. The level of DD depends on the assumptions; the timing does not.

\subsection{The verdict}
\label{sec:verdict}
This is the number I built the entire project to find out, so I want to state it plainly, the way I'd say it out loud in a viva, before I qualify it.

\begin{keybox}{Verdict, in my own words}
\textbf{The market-implied signal moved first, every single time.} In each of the three downgrade episodes I can date precisely, my computed distance-to-default series collapsed on the exact trading day the stress news hit — ahead of the first matching agency action by 3 days (the April 2019 loss, to ICRA's 3 May downgrade), 42 days (the slide through October--November 2019, to CARE's mid-November cut), and 68 days (the September 2018 CEO-term ruling, to the first ICRA/CARE downgrades on 28 November). Once I corrected the timeline, I also found that India Ratings alone acted three separate times across the period — so the agencies were not frozen, they were simply and consistently late relative to what the share price had already told anyone watching it.

\textbf{But neither signal timed the actual intervention.} On 4 March 2020 — the day before the moratorium — my own DD series read 3.45, its highest value since September 2018. I did not expect this result when I ran the code, and I want to be honest that it initially looked like a bug. It wasn't. The regulator acted on private supervisory information that neither the market nor the agencies had: the bank's December-2019 quarter, with a loss of roughly Rs 18,560 crore, was not published until 13 March 2020, eight days \emph{after} the moratorium. What my model gave was a long, correct, chronic warning that the bank was in the danger zone for the better part of eighteen months. What the ratings gave was a slower, stepwise confirmation of the same chronic story. Neither one predicted the specific week the regulator would act.
\end{keybox}

Two honest qualifications I want to make myself, before anyone reading this has to make them for me. First, the market and the agencies were often reacting to the \emph{same} public news at different speeds — the RBI's ruling on the CEO term, the first quarterly loss, the divergence disclosure. So when I say the market was ``faster,'' part of what I actually mean is ``a market re-prices every day, a committee meets every few months,'' not ``the market had private information the agencies lacked.'' Second, a chronically low DD is a mediocre early-warning alarm in practice: the distance also collapsed in September 2018 and then partially recovered by early 2019 without the bank actually failing, so if I had set a single fixed alarm threshold in advance, it would have fired a long time before the real event and stayed on for months — the kind of false-alarm fatigue that would make a risk desk start ignoring it. I think both qualifications are part of the honest answer, not asterisks to bury in a footnote.

\section{The limitations: why Merton is shaky for banks}
\label{sec:limits}
I don't want to present my own verdict above and then quietly move on, because the model that produced it has real cracks, and some of them are specific to the exact kind of company I pointed it at. Here is what I think I'd need to fix, or at least watch, before I'd trust this model unsupervised on another bank.
\begin{enumerate}
\item \textbf{The default point is ambiguous.} A bank's main liabilities are deposits. They are contractually short-term (many can be withdrawn on demand), so they behave like debt due now, yet they are not traded bonds with a face value and maturity date. Merton needs one debt amount due on one date. Choosing $D$ as total liabilities, as ``short-term plus half of long-term'', or as something else changes the level of $\DD$ materially (see the robustness table). For banks, failure usually comes as a \emph{run}: funding leaves before assets fall below liabilities. Merton has no run in it.
\item \textbf{Asset quality is opaque.} $V$ is the market value of loans nobody can see into. RBI divergence disclosures show that reported and supervisor-assessed non-performing assets differed by thousands of crores. The model's $V$ is inferred from the equity price, so it is only as good as the market's information.
\item \textbf{Regulatory forbearance can mask true values.} Rules that allow banks to delay recognising bad loans, restructure them or stay open while thinly capitalised keep reported assets above their economic value. A bank can be insolvent in economic terms and still be allowed to operate; the ``default'' event is then a regulatory decision, not a market outcome.
\item \textbf{The credit spread puzzle.} Empirically, Merton-type models tend to \emph{underestimate} real-world credit spreads for investment-grade issuers (for example Huang and Huang \cite{huang2012}); the normal distribution gives too little weight to sudden large losses. This is why KMV replaced $\Ncdf(-\DD)$ with an empirical EDF.
\item \textbf{Estimation error.} $V$, $\sigma_V$, $D$, $\mu$ and the volatility window are all chosen or estimated. As seen above, an old volatility spike leaving a window, or a high leverage ratio making the drift term dominate, can move DD by more than the underlying credit changed.
\end{enumerate}


\part{The Build: What I Actually Did, In Order}
\label{part:build}
\section{Why I built it in this order}

I want to be honest about something most reports hide: I did not write the Yes Bank pipeline first and trust it blindly. I built a small, fully-controlled toy problem first, where I already knew the right answer by hand, and I only pointed the machinery at a real bank once the toy problem came back clean. This section walks through that in the order I actually did it, with the real numbers my terminal printed, because I think the numbers mean more when you see them arrive the way I saw them arrive, rather than cleaned up and presented as if they always looked right.

\section{Step one: proving the simulator before I trust it}
\label{sec:build-mc}

My very first worry was simple: what if I had made an algebra mistake somewhere in Section~\ref{sec:dd}, and my distance-to-default formula was quietly wrong? A wrong formula will still produce a number — it just won't be the \emph{right} number — so I needed an independent check that didn't go through the same algebra a second time. That's what Section~\ref{sec:mc}'s Monte Carlo simulation is for, and here is exactly what it printed the first time I ran it, for the illustrative firm with $V_0=120$, $D=100$, $\mu=8\%$, $\sigma=25\%$, $T=1$:

\begin{verbatim}
DD = 0.9243    closed-form PD = N(-DD) = 17.7669%

    paths   MC PD (V_T<D)   +/- 95% CI  abs error
      100        13.0000%      6.5915%     4.7669%
    1,000        18.8000%      2.4217%     1.0331%
   10,000        18.1000%      0.7546%     0.3331%
  100,000        17.8270%      0.2372%     0.0601%
1,000,000        17.7089%      0.0748%     0.0580%

Barrier-any-time PD (100k paths, 252 steps): 37.7270%  (always >= terminal PD)
\end{verbatim}

\begin{defbox}{reading this output the way I had to read it the first time}
The top line is the answer my formula \emph{claims}: a 17.77\% chance the firm's assets end below its debt in one year. Everything below it is me checking that claim by literally simulating the firm's future a million separate times and counting. At 100 random futures, my estimate (13.00\%) is nowhere near the formula's answer — but the ``$\pm$95\% CI'' column already warned me that with only 100 trials, I should expect to be off by as much as 6.59 percentage points just from randomness, so this isn't a failure, it's exactly what a small sample should look like. As I scaled up to a million simulated futures, my estimate (17.7089\%) landed within 0.058 percentage points of the formula's 17.7669\% — well inside the confidence band the statistics predicted. That convergence is what told me my algebra in Section~\ref{sec:dd} was actually correct, not just plausible-looking.

The barrier line is a separate, harder question I asked out of curiosity: what if I count a default any time the simulated path dips below $D$ during the year, not just at the one-year mark? That number came back at 37.73\% — more than double the terminal-only answer. This makes sense once I thought about it: a path has many more chances to dip below a line at some point during a whole year than it has chances to specifically end below that line on the very last day. It's the difference between asking ``will the river ever flood this year'' versus ``will the river be flooded on exactly December 31st'' — the first question is easier to answer yes to. Merton's original model only asks the second, narrower question, and that is itself a limitation I come back to in Section~\ref{sec:limits}.
\end{defbox}

Once I had this convergence in hand, I trusted the formula enough to point it at a real company.

\section{Step two: turning the formula into a daily pipeline}
\label{sec:build-pipeline}

The toy problem has one $V_0$, one $\sigma$, one answer. Yes Bank has 573 days of prices, a balance sheet that changes every quarter, and a question that only makes sense as a \emph{time series}: what did distance-to-default look like on every single trading day, not just once? This meant I had to build \texttt{yes\_bank\_dd.py} to do, automatically, for every day in my CSV, the same two-equation solve I describe by hand in Section~\ref{sec:estimate} — pull that day's equity value and the trailing 60-day equity volatility, feed them in alongside the most recently published balance sheet, and solve for that day's $V$, $\sigma_V$, $\DD$, and $d_2$.

Here is a genuine excerpt of what that script printed, taken straight from my own terminal, one row per month-end for readability (the real file has one row per trading day):

\begin{verbatim}
Date         Price   sigma_E       V        sigma_V    DD      d2
2018-04-30  362.00   0.356   308260.415   0.097   3.645   3.230
2018-05-31  346.20   0.337   349962.452   0.077   3.847   3.328
2018-06-30  339.65   0.294   348451.464   0.066   4.423   3.819
2018-07-31  367.95   0.229   371937.976   0.052   5.697   4.932
2018-08-31  343.50   0.284   366289.863   0.062   4.588   3.938
2018-09-30  183.65   0.812   326683.137   0.121   1.342   1.012
2018-10-31  188.10   0.905   362699.313   0.132   1.083   0.780
2018-11-30  169.80   1.000   355724.074   0.145   0.843   0.567
2018-12-31  181.80   0.679   365225.021   0.084   1.844   1.370
2019-01-31  194.10   0.619   373225.178   0.078   2.086   1.574
2019-02-28  231.15   0.712   380817.321   0.109   1.660   1.292
2019-03-31  275.10   0.685   391170.631   0.119   1.732   1.396
2019-04-30  168.00   0.967   365584.151   0.132   0.936   0.633
2019-05-31  147.80   0.855   364292.321   0.096   1.293   0.877
2019-06-30  108.75   0.928   354566.361   0.084   1.165   0.690
2019-07-31   91.20   0.824   343842.953   0.060   1.591   0.927
2019-08-31   59.95   0.938   337125.300   0.056   1.364   0.649
2019-09-30   41.40   1.010   332214.115   0.046   1.360   0.485
2019-10-31   70.40   1.351   322672.404   0.146   0.159  -0.114
2019-11-30   68.30   1.250   304734.586   0.125   0.378   0.057
2019-12-31   46.95   1.104   307080.227   0.067   0.905   0.305
2020-01-31   39.25   0.726   309748.697   0.027   2.668   1.163
2020-02-29   34.60   0.711   308663.091   0.023   2.956   1.204
2020-03-31   16.15   1.844   270583.950   0.126  -0.571  -0.889
\end{verbatim}

I want to actually walk through this table rather than just show it, because every column here has a story, and I learned the model far better by reading my own output line by line than I did by deriving it on paper.

\paragraph{July 2018 — the safest Yes Bank ever looked in my data.} $\DD=5.697$, the highest value in the whole series. This is the month CARE gave Yes Bank its AAA upgrade, and my model agrees that, by its own lights, the bank genuinely looked rock-solid at that moment: equity volatility was a low 22.9\%, and the implied asset cushion over the default point was enormous. I want to flag something important here, though — my model agreeing with CARE in July 2018 is not the model being vindicated. Both were reading the same calm market and the same reported numbers. The test of the model isn't whether it agrees in the good months; it's what happens next.

\paragraph{September 2018 — the number that moves in one line.} Between August 31 ($\DD=4.588$) and September 30 ($\DD=1.342$), distance to default fell by more than three full standard deviations in a single calendar month. This is the CEO-term shock hitting the share price. Notice what drove it mechanically: $\sigma_E$ nearly triples, from 28.4\% to 81.2\%, in the same window. A sudden crash doesn't just lower the price — it massively raises the \emph{measured volatility}, because my 60-day window now contains that huge one-day move. Both effects — lower $V$ (since $V$ is solved from a lower $E$) and higher $\sigma_V$ (since $\sigma_V$ is solved from a higher $\sigma_E$) — push $\DD$ down at the same time, through the denominator and the numerator together. This is exactly the mechanism behind the sharp drop I describe in Section~\ref{sec:verdict}, except now I'm looking at the actual two numbers that caused it, not just the headline result.

\paragraph{Why $V$ barely falls even as the share price collapses.} This genuinely confused me the first time I looked at this table, so I want to walk through the confusion honestly. Between July 2018 ($V=371{,}938$) and September 2019 ($V=332{,}214$), Yes Bank's \emph{share price} fell from Rs 368 to Rs 41 — a 89\% collapse. But my solved \emph{asset value} $V$ only fell from about Rs 3.72 lakh crore to about Rs 3.32 lakh crore — about 11\%. At first I thought this was a bug in my solver. It isn't; it's the model telling me something real about leverage. Yes Bank's equity ($E$, the market cap) was always a small sliver on top of a much larger pile of liabilities ($D$, in the hundreds of thousands of crore). When a firm is this leveraged, a given percentage fall in the small equity sliver corresponds to a much smaller percentage fall in the (much larger) total asset value underneath it — the same rupee loss is a huge percentage of a small number and a tiny percentage of a big number. This is actually the mirror image of why $\sigma_E$ (92\% in some months) is so much larger than $\sigma_V$ (4.6\% to 14.6\% throughout): leverage amplifies the equity holder's percentage volatility relative to the firm's underlying asset volatility, in exactly the way equation~\eqref{eq:e2} says it should. Seeing this in my own numbers, rather than just in the symbols, is what made equation~\eqref{eq:e2} finally feel real to me instead of just algebraically true.

\paragraph{October 2019 — the one month $d_2$ actually went negative.} $\DD=0.159$, $d_2=-0.114$. A negative $d_2$ means the risk-neutral default probability $\Ncdf(-d_2)$ was \emph{above} 50\% for that reading — the risk-neutral world, which I built in Section~\ref{sec:merton} specifically to price debt rather than to predict outcomes, was by this point pricing Yes Bank's debt as more likely to default than not. I want to be precise about what this does and doesn't mean: it does not mean I, or the model, am claiming Yes Bank was actually more likely than not to fail that month. It means that if you wanted to replicate Yes Bank's bonds using a risk-free asset and the stock, the maths implied a default probability above a coin-flip under the pricing measure. That's a market-implied pricing statement, not a forecast — exactly the distinction I built in Section~\ref{sec:physical-rn}, and this is the one month in my whole dataset where the distinction stopped being academic and started being the headline number itself.

\paragraph{The March 2020 whiplash, in the actual numbers.} $\DD$ climbs from 0.905 (Dec 2019) to 2.668 (Jan 2020) to 2.956 (Feb 2020) — recovering most of the way back toward its 2018 level — before crashing to $-0.571$ by the 31 March close, after the moratorium had already happened. I sat with this for a while before writing Section~\ref{sec:verdict}'s verdict, because my honest first reaction was that the model had failed at the one moment that mattered most. What I now think happened, mechanically, is visible right there in the table: $\sigma_E$ drops to 0.027 (2.7\%) in February 2020 — an unusually \emph{calm} reading, because the huge volatility spike from October 2019 had, by February, aged out of my 60-day trailing window, while the actual news driving the bank toward the cliff (an unpublished quarterly loss) hadn't happened yet in a way the market could price. The model wasn't lying to me in February 2020. It was accurately reporting that the \emph{priced, observable} risk had calmed down, right before a shock that nobody outside the RBI could see arrived from outside the market entirely. That is, I think, the single most important sentence in this whole project, and I did not understand it until I'd built the pipeline and watched it produce a number I didn't expect.

\section{Building the repository}

Once the numbers were right, I didn't want this to live only as a PDF. I built a proper repository: a \texttt{src/} folder holding \texttt{yes\_bank\_dd.py} and \texttt{merton\_monte\_carlo.py}, a \texttt{data/} folder for the raw price CSV, an \texttt{outputs/} folder that the pipeline writes its computed series into, a \texttt{figures/} folder for every chart in this note, a \texttt{slides/} folder for the presentation version of this argument, and a \texttt{docs/} folder holding this note itself and its \LaTeX{} source. I wrote a README that states the headline verdict in the first paragraph, because I didn't want anyone — a recruiter, a professor, future-me — to have to read twenty pages to find out what I actually found. I licensed the code MIT, added a \texttt{.gitignore} and a \texttt{requirements.txt}, and made the price-lookup logic in \texttt{yes\_bank\_dd.py} search for any CSV file whose name contains ``yes'' and ``bank'' in its own folder or a sibling \texttt{data/} folder, rather than hard-coding one exact filename — a small thing, but it's the difference between a script that only ever runs on my laptop and one a stranger could actually clone and run.


\part{What This Means in the Real World}
\label{part:realworld}
\section{So what? Why any of this should matter to someone who isn't grading it}

Everything up to this point has been building a model and testing it. This section is me stepping back and asking the question I'd actually be asked by a risk manager, an investor, or an interviewer: fine, you built a working distance-to-default pipeline and it predicted Yes Bank's stress eighteen months early — so what should anyone actually \emph{do} differently because of this? I think there are real, specific answers, for several different readers, and I want to give each of them their own answer rather than one vague conclusion.

\subsection{For a bank investor or a bond holder}
If I had been holding Yes Bank's Tier~II bonds or shares in mid-2018, my model's own numbers say I had a genuinely usable early-warning signal available to me, for free, from public price data, starting in September 2018 — more than a year before the first rating agency action that seriously questioned Yes Bank's capital position. The honest caveat, which I don't want to undersell, is in my own robustness check: a DD-based alarm would also have fired falsely, in the sense that Yes Bank survived (as an independent-ish entity with new management) for over a year after the alarm first went off, and partially recovered in early 2019 without failing. A real investor using this signal in real time would have needed to decide, in September 2018, whether a four-standard-deviation collapse in one month was a reason to exit entirely or a reason to just reduce exposure and watch closely — and my model, by itself, doesn't tell you which. What it does tell you is that waiting for a rating agency to confirm the story before acting means acting, on average across my three measured episodes, somewhere between three days and over two months later than the market had already told you.

\subsection{For a bank's own risk management desk}
This is where I think the project's biggest practical lesson sits, and it isn't ``use Merton instead of ratings.'' It's that the two signals fail in \emph{opposite, complementary} ways, which is exactly why a real risk desk should want both running simultaneously, not one or the other. My DD series is loud, immediate, and chronically pessimistic once a firm is under real stress — exactly the kind of signal you want for day-to-day monitoring, precisely because it updates every single trading session and costs nothing to compute beyond a share price. But it is structurally blind to anything that happens outside the market's view: it could not see the December-2019 loss before 13 March 2020 because the market itself couldn't see it. A rating agency, for all its slowness, has analysts who can sit across a table from management, ask for internal numbers the market never sees, and — this is the part I think gets underrated in most critiques of agencies, including the one I opened this whole project with — occasionally catch something the market hasn't priced in yet \emph{because it isn't public}. A risk desk that only watched DD would have missed nothing about Yes Bank's ongoing distress but would have been just as surprised by the specific timing of the moratorium as everyone else. A risk desk that only watched ratings would have been months late to the party entirely. Running both, and specifically watching for the moments they \emph{disagree}, is the actual operational lesson I'd pull from this project if I were building a monitoring dashboard rather than writing a note about one.

\subsection{For a regulator, and for Indian financial regulation specifically}
I think my results are an argument, even if a small, single-case one, for why the RBI increasingly leans on its own direct, non-public supervisory channels — on-site inspections, divergence audits, confidential reporting — rather than relying on either market prices or published ratings as its primary early-warning tool. Both of the public signals I tested were, in their own ways, informative but insufficient: the market was fast but blind to what it couldn't see, and the agencies were seeing much of the same public information the market saw, just slower. The one genuinely decisive, non-public piece of information in this entire case — Yes Bank's real December-2019 numbers — only became visible to the RBI through its own supervisory access, not through anything a market-based model or a public rating could have picked up in time. If I were advising a regulator rather than a student writing a project, my takeaway wouldn't be ``trust market signals more.'' It would be ``market signals are a genuinely useful early, cheap, continuously-updating tripwire for \emph{where} to point your scarce supervisory attention — they are not a substitute for having supervisory access in the first place.''

\subsection{The specific, structural problem with the rating agencies, now that I've actually checked it}
When I started this project, my working assumption — shaped partly by the general global critique in Section~\ref{sec:critique} and partly, I'll admit, by something closer to received wisdom than verified fact — was that the rating agencies had probably been pressured into going slow on Yes Bank specifically. Having now actually read SEBI's 2023 order, I think the real, verified picture is more useful than the simpler story I started with, and it's worth spelling out why. SEBI's forensic audit \emph{did} find that CARE's senior management interfered with and delayed a rating action — for DHFL, a different, related stressed borrower, in the same period. For Yes Bank specifically, the same investigation, after years of cross-examination, concluded the interference charge couldn't be sustained on the evidence available. I think the honest, structural lesson here is not ``CARE was corrupt about Yes Bank'' (the regulator looked hard and didn't find enough to prove that) but something a little more unsettling precisely because it's more mundane: \textbf{the same issuer-pays conflict of interest that was proven to have delayed one bank's downgrade was structurally present, and plausible, across the whole stressed-lender portfolio of that period} — DHFL, IL\&FS, and Yes Bank were all being rated by the same analysts, under the same commercial pressure, in the same eighteen-month window, and the one case SEBI could actually prove is not obviously a special case rather than a sample. A model like mine doesn't have this problem at all, for a very simple, almost dull reason: a share price is set by thousands of independent, anonymous traders with money at risk, not by a committee whose employer is paid by the entity being graded. That structural independence — not any claim about speed or sophistication — is, I now think, the real argument for keeping a market-based signal running alongside agency ratings, and it's a stronger and more specific argument than the one I'd have made before I did the verification work in Section~\ref{sec:case}.

\subsection{What I'd tell my own FRM-prep self, a year ago}
I'm building this project partly as an FRM candidate, so I want to end this section on the question that's most practical for me personally: what does this change about how I'd answer a credit-risk question on an exam, or in an interview? Three things, concretely. First, I now actually understand \emph{why} distance-to-default is quoted in standard-deviation units rather than rupees or percentages — because I had to solve the two-equation system myself and watch what happened to $\sigma_V$ as leverage changed, rather than just memorising the formula. Second, I now have a real, specific, defensible answer to ``what's wrong with structural models for banks'' that goes beyond the textbook line about deposits being hard to model — I have an actual month (October 2019) where I watched $d_2$ go negative in my own data and had to sit with what that number was and wasn't telling me. Third, and this is the one I didn't expect going in: I now think the single most exam-relevant, interview-relevant fact from this whole project is the March 2020 whiplash itself — a model can be \emph{correctly} reporting that priced risk has fallen while the \emph{true} risk is rising, because the two things aren't the same, and the gap between them is exactly where supervisory information earns its keep. That's a sentence I could only have written after watching my own pipeline produce a number that surprised me.


\part{Closing the Loop}
\label{part:close}
\section{The synthesis and close}
\label{sec:synth}
I'll say this plainly, because I think it's the single sentence this entire project was built to earn the right to say: the two approaches fail in opposite ways, and neither one is optional. Ratings are slow, stepwise and dependent on disclosure, but they can weigh private information, management quality and regulatory context, and they anchor mandates and covenants in a way a raw number from a model never could. Market-implied models are continuous, fast and cheap to update, but they are noisy, blind to losses the market cannot see, and structurally awkward for banks specifically, for the reasons I set out in Section~\ref{sec:limits}. Yes Bank shows both faults in the same case study: the market signalled danger more than a year before the agencies caught up, and neither the market nor the agencies anticipated the specific week of the RBI's intervention.

\begin{keybox}{Take-away, in my own words}
I didn't go looking for a tidy answer, and I don't think I found one — I found a real, specific, dated answer instead, which I think is more useful. Neither agency ratings nor market-implied models are sufficient alone. A bank, an investor, or a regulator that only watches one of them is choosing which kind of blind spot to live with. Real risk management uses both: the committee's slower judgement for what can't be priced from the outside, and the market's faster price for what changes every single day whether a committee is in session or not.
\end{keybox}

Where this idea led next, as a pointer rather than core content: Vasicek's one-factor extension of the Merton model underlies the Basel IRB regulatory capital formula — the same equity-as-option logic I built in Part~\ref{part:theory}, scaled up into the formula banks use to compute regulatory capital against an entire loan portfolio rather than a single firm.

\section{What I verified, what changed, and what I still could not confirm}
\label{sec:verify}
I want this section to be an honest accounting, not a formality. When I first wrote this note, several dates below were marked as unverified. I went back and checked each one against a primary source — a company filing, an agency's own release, or in one case an actual SEBI adjudication order — and most of them are now resolved. I'm listing what changed as well as what's still open, because I think showing the correction is more useful to a reader than silently fixing it and pretending the first draft never happened.

\begin{itemize}
\item \textbf{Resolved --- the RBI's letter on the CEO's term.} I now have this to a primary source: the RBI's letter is dated 17 September 2018; Yes Bank disclosed it to the exchanges on 19 September; the sharpest part of the market reaction landed on 21 September \cite{bsSep19disc,bsSep21fall}.
\item \textbf{Resolved --- CRISIL's pre-crisis silence.} I confirmed this rather than just assuming it: I found no CRISIL rating action on Yes Bank's debt before the 19 March 2020 CD rating. ICRA, CARE, India Ratings and Moody's were the active agencies throughout the stress period.
\item \textbf{Resolved, and corrected --- the CARE/SEBI interference finding.} I initially had this wrong in a way I think is worth being transparent about: I wrote down that SEBI had penalised CARE specifically for delaying its Yes Bank rating. Having read SEBI's actual 2023 order \cite{sebiCareOrder}, the finding is more specific: interference was proven for CARE's handling of DHFL; for Yes Bank specifically, the same order states the interference charge could not be sustained on the evidence. I've corrected the claim in Section~\ref{sec:case} rather than leaving the stronger, inaccurate version in place.
\item \textbf{Resolved --- intermediate rating steps I was missing.} My first timeline had a large, unexplained gap between CARE's November 2018 downgrade and its November 2019 cut. On re-verification I found three India Ratings downgrades I'd missed entirely (8 May 2019, 31 August 2019, 19 December 2019) plus a fourth on 12 February 2020 — meaning the agencies were considerably more active across 2019 than my first draft gave them credit for, even though each individual action still lagged my computed DD series.
\item \textbf{Resolved --- CARE's October 2018 watch date.} Confirmed to CARE's own press release, reported the same day: 1 October 2018 \cite{bsCareOct1}.
\item \textbf{Resolved --- India Ratings' May 2019 action.} Confirmed: the downgrade was filed late on 8 May 2019, reported the following morning \cite{bsIndraMay19}.
\item \textbf{Resolved --- SBI's stake and the lock-in.} Confirmed against the Reconstruction Scheme's own text \cite{yesReconScheme} and SBI's own disclosures \cite{sbiStakeScheme}: the Scheme capped the investor bank's stake at 49\% and required a minimum of 26\% for three years from infusion; SBI's initial tranche gave it roughly 48.2\%, building toward a total commitment of Rs 7,250 crore. The three-year lock-in genuinely did run until 13 March 2023, which I can confirm independently since Yes Bank's stock fell when that lock-in expired on schedule.
\item \textbf{Still open --- a few balance-sheet interpolations.} I still had to interpolate the ``other liabilities'' line for June 2018 and December 2018, and the whole December 2017 figure, since I could not find a disclosed breakdown for those exact dates. I also still assume roughly 30-day publication lags for a few quarters where I used the actual filing date for others. These are my modelling choices, not facts I failed to look up, and I flag them as such rather than as things still to verify.
\item \textbf{Still open --- the exact QIP date.} Yes Bank's own press release names the month (August 2019) but not the exact day; I used 15 August as a working assumption.
\item \textbf{Not primary-sourced, and not meant to be --- Enron and the 2007--09 crisis.} These remain general historical background in Section~\ref{sec:critique}, included for context, not as claims I'm treating as load-bearing for my own result.
\item \textbf{Expected to vary --- Monte Carlo output.} The exact figures in Section~\ref{sec:build-mc} will shift slightly with a different random seed or machine; the convergence pattern will not.
\end{itemize}

\appendix
\section{Glossary}
\begingroup\small
\begin{description}[leftmargin=2.4cm,labelwidth=2.2cm,itemsep=1pt]
\item[Asset volatility] Yearly standard deviation of a firm's asset returns ($\sigma_V$); not directly observable.
\item[Call option] The right to buy an asset at strike $K$ at expiry; pays $\max(S_T-K,0)$.
\item[CAGR] Compound annual growth rate: the constant yearly growth that turns a start value into the end value; a geometric average.
\item[Credit spread] Extra yield a risky bond pays over the risk-free rate.
\item[Default point] The debt level the assets must cover ($D$); the model's threshold for default.
\item[Distance to default] Number of standard deviations between expected asset value at the horizon and the default point.
\item[EDF] Expected default frequency: KMV's empirical probability, read from its history of defaults for firms with a given DD.
\item[GBM] Geometric Brownian motion: a random path with random percentage steps.
\item[Lognormal] A variable whose logarithm is normally distributed; always positive.
\item[Moratorium] A legal suspension of a bank's obligations (here, capped withdrawals) under Section 45 of the Banking Regulation Act, 1949.
\item[Monte Carlo simulation] Estimating a probability by generating many random scenarios and counting outcomes, rather than solving a formula.
\item[Reconstruction Scheme] The RBI/government order under which a failed bank is restructured and recapitalised by new investors; for Yes Bank, notified 13 March 2020.
\item[Risk-neutral] Priced as if investors demanded no premium for risk; all assets grow at $r$.
\item[Structural model] A credit model that derives default from assets versus debt.
\item[Through-the-cycle] A rating philosophy that aims to hold across economic ups and downs.
\item[Volatility drag] The reduction in compound growth caused by volatility, about $\tfrac12\sigma^2$ per year.
\end{description}
\endgroup

\begin{thebibliography}{99}\small
\bibitem{merton1974} R.\,C.\ Merton (1974). On the pricing of corporate debt: the risk structure of interest rates. \emph{Journal of Finance} 29(2), 449--470.
\bibitem{blackScholes1973} F.\ Black and M.\ Scholes (1973). The pricing of options and corporate liabilities. \emph{Journal of Political Economy} 81(3), 637--654.
\bibitem{blackcox1976} F.\ Black and J.\,C.\ Cox (1976). Valuing corporate securities: some effects of bond indenture provisions. \emph{Journal of Finance} 31(2), 351--367.
\bibitem{crosbie2003} P.\ Crosbie and J.\ Bohn (2003). \emph{Modeling Default Risk}. Moody's KMV.
\bibitem{huang2012} J.-Z.\ Huang and M.\ Huang (2012). How much of the corporate-treasury yield spread is due to credit risk? \emph{Review of Asset Pricing Studies} 2(2), 153--202.
\bibitem{hull} J.\ Hull. \emph{Options, Futures, and Other Derivatives} (chapters on Black--Scholes and credit risk).
\bibitem{malz2011} A.\ Malz (2011). \emph{Financial Risk Management: Models, History, and Institutions}. Wiley.
\bibitem{bsCareWatch} Business Standard (1 Oct 2018). CARE Ratings put certain YES Bank debt instruments on credit watch. \url{https://www.business-standard.com/article/finance/care-ratings-put-certain-yes-bank-debt-instruments-on-credit-watch-118100100144_1.html}
\bibitem{bsIcraWatch} Business Standard (17 Nov 2018). Icra places Yes Bank's rating on watch with negative effect over RBI order. \url{https://www.business-standard.com/amp/article/finance/icra-places-yes-bank-s-rating-on-watch-with-negative-effect-over-rbi-order-118111700418_1.html}
\bibitem{bsNov28} Business Standard/Reuters (29 Nov 2018). Yes Bank hits over two-year low after ICRA, CARE cut ratings. \url{https://www.business-standard.com/amp/article/reuters/yes-bank-hits-over-two-year-low-after-icra-care-cut-ratings-118112900183_1.html}; and Business Standard (29 Nov 2018), Yes Bank hits 52-week low after ratings downgrade (Moody's).
\bibitem{bsQ3FY19} Business Standard (24 Jan 2019). Yes Bank vaults on high volume after Q3 result. \url{https://www.business-standard.com/article/news-cm/yes-bank-vaults-on-high-volume-after-q3-result-119012400681_1.html}
\bibitem{bsDivFY16} Business Standard (12 May 2017). Yes Bank under-reported NPAs by over Rs 4,000 crore in FY16. \url{https://www.business-standard.com/article/pti-stories/yes-bank-under-reported-npas-by-over-rs-4-000-crore-in-fy16-117051201439_1.html}
\bibitem{bsDivFY18} Business Standard (13 Feb 2019). RBI finds no divergence in provisioning, asset classification: Yes Bank. \url{https://www.business-standard.com/amp/article/pti-stories/rbi-finds-no-divergence-in-provisioning-asset-classification-yes-bank-119021301416_1.html}
\bibitem{bsIcraMay19} Business Standard (4 May 2019). Icra downgrades YES Bank's bonds as quality of large borrowers dwindles. \url{https://www.business-standard.com/amp/article/finance/icra-downgrades-yes-bank-s-bonds-as-quality-of-large-borrowers-dwindles-119050400299_1.html}
\bibitem{bsFall19} Business Standard (6 May 2019). YES Bank falls 5\% as ICRA downgrades bonds rating. \url{https://www.business-standard.com/amp/article/markets/yes-bank-falls-5-as-icra-downgrades-bonds-rating-119050600162_1.html}
\bibitem{bsIndraMay19} Business Standard (9 May 2019). India Ratings downgrades Yes Bank's long-term ratings. \url{https://www.business-standard.com/amp/article/news-ani/india-ratings-downgrades-yes-bank-s-long-term-ratings-119050900352_1.html}
\bibitem{icraJul19} ICRA (24 Jul 2019). Yes Bank Limited: Long-term ratings downgraded (rationale). \url{https://www.icra.in/Rating/GetRationalReportFilePdf/82233~Yes\%20Bank_r_24072019.pdf}
\bibitem{bsIcraJul19} Business Standard (25 Jul 2019). ICRA downgrades YES Bank's long-term rating for bonds worth Rs 32,911 crore. \url{https://www.business-standard.com/amp/article/economy-policy/icra-downgrades-yes-bank-s-long-term-rating-for-bonds-worth-rs-32-911-crore-119072500053_1.html}
\bibitem{ycQ2} YES BANK (1 Nov 2019). Press release: financial results for the quarter ended 30 September 2019. \url{https://www.yes.bank.in} (investor relations, Q2 FY20 press release).
\bibitem{bsCareNov19} Business Standard (14 Nov 2019). Yes Bank gets revision in credit ratings from CARE. \url{https://www.business-standard.com/article/news-cm/yes-bank-gets-revision-in-credit-ratings-from-care-119111400471_1.html}
\bibitem{bsDiv19} Business Standard (20 Nov 2019). RBI found Rs 3,277-cr divergence in NPA provisioning for FY19: YES Bank. \url{https://www.business-standard.com/amp/article/finance/rbi-found-rs-3-277-cr-divergence-in-npa-provisioning-for-fy19-yes-bank-119111901799_1.html}
\bibitem{bsIndraDec19} Business Standard (19 Dec 2019). Yes Bank receives downward revision in ratings; placed on rating watch negative. \url{https://www.business-standard.com/amp/article/news-cm/yes-bank-receives-downward-revision-in-ratings-placed-on-rating-watch-negative-119121900769_1.html}
\bibitem{bsCareDec19} Business Standard (1 Jan 2020). Yes Bank slips on rating downgrade. \url{https://www.business-standard.com/amp/article/news-cm/yes-bank-slips-on-rating-downgrade-120010100165_1.html}
\bibitem{icraMar20} ICRA (6 Mar 2020). Yes Bank Limited: Ratings downgraded. \url{https://www.axistrustee.in/pdf/Yes\%20Bank\%20Limited\%20-(ICRA)\%206th\%20March\%202020.pdf}
\bibitem{bsMar9} Business Standard (9 Mar 2020). Rating agencies downgrade Yes Bank. \url{https://www.business-standard.com/article/news-cm/rating-agencies-downgrade-yes-bank-120030900320_1.html}
\bibitem{btQ3} Business Today (15 Mar 2020). YES Bank Q3 net loss spikes to Rs 18,564 cr, its worst ever. \url{https://www.businesstoday.in/latest/corporate/story/yes-bank-suffers-rs-9000-crore-in-december-quarter-252050-2020-03-14}
\bibitem{crisilMar20} CRISIL (19 Mar 2020). YES Bank Limited: rating rationale. \url{https://www.crisil.com/mnt/winshare/Ratings/RatingList/RatingDocs/YES_Bank_Limited_March_19_2020_RR.html}
\bibitem{angel3q19} Angel Broking. Yes Bank 1QFY2019 and 3QFY2019 result updates (quarterly deposits and borrowings).
\bibitem{screenerYes} Screener.in. Yes Bank Ltd, consolidated balance sheet (March 2019). \url{https://www.screener.in/company/532648/consolidated}
\bibitem{bsSep19disc} Business Standard / IANS (19 Sep 2018). Yes Bank CEO Rana Kapoor's tenure to end on Jan 31, 2019. \url{https://www.business-standard.com/article/news-ians/yes-bank-ceo-rana-kapoor-s-tenure-to-end-on-jan-31-2019-118091901101_1.html}
\bibitem{bsSep21fall} Business Standard / PTI (21 Sep 2018). Yes Bank sinks 34\%; m-cap erodes by Rs 14,452 cr as RBI trims Rana Kapoor's term as CEO. \url{https://www.business-standard.com/article/pti-stories/yes-bank-sinks-34-m-cap-erodes-by-rs-14-452-cr-as-rbi-trims-rana-kapoor-s-term-as-ceo-118092100272_1.html}
\bibitem{bsCareOct1} Business Standard (1 Oct 2018). Asset quality outlook `stable', proportion of NPAs down in a year: YES Bank. \url{https://www.business-standard.com/article/finance/asset-quality-outlook-stable-proportion-of-npas-down-in-a-year-yes-bank-118100100416_1.html}
\bibitem{bsIndraAug19} Business Standard / PTI (31 Aug 2019). India Ratings downgrades YES Bank to `IND A+', says `outlook is negative'. \url{https://www.business-standard.com/amp/article/pti-stories/india-ratings-downgrades-yes-bank-outlook-negative-119083100463_1.html}
\bibitem{bsIndraFeb20} Business Standard (12 Feb 2020). India Ratings downgrades Yes Bank issuer rating from `A' to `A-'. \url{https://www.business-standard.com/amp/article/finance/india-ratings-downgrades-yes-bank-issuer-rating-from-a-to-a-120021201750_1.html}
\bibitem{yesReconScheme} Reserve Bank of India / Government of India (13 Mar 2020). Yes Bank Limited Reconstruction Scheme, 2020, issued under Section 45(4)--(7) of the Banking Regulation Act, 1949.
\bibitem{sbiStakeScheme} Business Standard (6 Mar 2020 and 18 Mar 2020). YES Bank `won't fall off the cliff', SBI to rescue it with 49\% stake; SBI announces Rs 7,250-crore fund infusion into Yes Bank. \url{https://www.business-standard.com/amp/article/companies/sbi-to-rescue-yes-bank-with-49-stake-govt-says-depositors-money-safe-120030601650_1.html}
\bibitem{sebiCareOrder} Securities and Exchange Board of India (2023). Order in the matter of CARE Ratings Ltd. (WTM/ASB/DDHS/DDHS-RAC-POD-II/25845/2023-24). \url{https://www.sebi.gov.in/sebi_data/attachdocs/apr-2023/1681988405288_3.pdf}
\end{thebibliography}


\end{document}
